\section{Discusi\'on de resultados}

\subsection{Interpretacion de la comparacion mononodo vs. multimodo}

El cluster distribuido funciono correctamente y ejecuto el proyecto con estado final \texttt{SUCCEEDED}. Sin embargo, para este dataset el entorno mononodo fue mas rapido. La razon principal es que el volumen de datos es pequeno para un cluster: 58 MB en CSV y 222\,874 registros. En ese rango, el costo fijo de iniciar aplicaciones YARN, coordinar contenedores, leer desde HDFS, serializar datos, mover bloques de shuffle y sincronizar resultados puede superar el beneficio de paralelizar.

Tambien influyo el contexto de virtualizacion. Los nodos workers tienen recursos ajustados, especialmente 2 GB de RAM y 1 vcore cada uno. Durante algunas ejecuciones aparecieron advertencias \texttt{FetchFailed}; Spark pudo recuperarse mediante reintentos y la aplicacion finalizo correctamente, pero esos reintentos aumentan el tiempo total. En un dataset mayor, o con workers con mas memoria y CPU, el cluster deberia mostrar mejor escalamiento.

\subsection{Interpretacion de las consultas SQL}

La consulta por periodo muestra que el numero de emisiones aumento de 69\,880 en 2023 a 79\,013 en 2025. Serenazgo es el concepto con mayor monto total en los tres periodos, superando los 24 millones en 2023 y llegando a 27.8 millones en 2025. Barrido de calles tiene el monto total mas bajo, pero crece de 2.64 millones en 2023 a 3.31 millones en 2025.

La consulta de top 20 para 2025 permite identificar predios exclusivos con montos netos excepcionalmente altos. Estos registros son utiles para auditoria, priorizacion de fiscalizacion o analisis de concentracion de la recaudacion.

\subsection{Interpretacion de los modelos}

En regresion, el MLPRegressor logra el menor error global. El SVRegressor tiene menor MAE local, pero mayor RMSE y MSE, lo que sugiere mayor sensibilidad ante errores grandes. Los valores de $R^2$ moderados indican que la variable \texttt{MONTO\_SERENAZGO} depende de factores adicionales que no estan completamente representados en el dataset.

En clasificacion, Random Forest se comporta mejor que el MLPClassifier. Esto es razonable porque los arboles capturan relaciones no lineales y reglas por umbral con menor necesidad de ajuste fino. El MLPClassifier requiere mayor calibracion de arquitectura, iteraciones y posiblemente balanceo de clases para mejorar su F1-score.
